\documentclass[11pt]{article}
\usepackage{amsmath,amssymb}
\title{Room 12 follow-up: the continuity-tracking question, resolved}
\date{2026-09-27}
\begin{document}
\maketitle

\section*{The question}
Reviewer AW raised a serious concern (see \texttt{room12\_summary.tex}): \texttt{ratio\_amp}
computes $s=\sqrt{c^N}$ via an independent principal-branch call at each sampled point, with no
continuity tracking across nearby $N$. If the physically meaningful branch should vary
\emph{continuously} with $(N,a)$ (tracking the actual saddle asymptotic), then the observed
erratic/reciprocal-bimodal behaviour in Room 7's even-$a$ data might be substantially a branch-cut
discontinuity artifact of the tool, not a real feature of $A_+/A_-$.

\section*{The check}
Direct computation (Python/mpmath, 150 digits -- the genuinely huge dynamic range here,
$|c^N|\sim10^{-96}$ at these $(N,a)$, needs this much precision; earlier attempts at 50--60 digits
returned false ``degenerate'' flags from an over-strict, non-relative threshold, not a real
computational failure), $q=29$, $N=277$ to $303$ even-$a$ (12 points), comparing:
\begin{itemize}
\item \textbf{Principal branch}: $s=\sqrt{c^N}$, mpmath's own principal square root, independently
at each $N$ (what \texttt{ratio\_amp} does).
\item \textbf{Continuity-tracked branch}: at each $N$, pick whichever of $\pm\sqrt{c^N}$ is closer
to the previous point's chosen value.
\end{itemize}

\textbf{Result: the continuity-tracked sequence STILL oscillates between the same two clusters}
($\approx0.013$--$0.015$ and $\approx65$--$78$) at essentially every step, exactly as the naive
principal-branch sequence does. \textbf{The reciprocal-bimodal/erratic phenomenon is NOT a branch-
tracking artifact; it is a genuine feature of $A_+/A_-$ as $N$ varies at fixed seed $a/29$.}

\textbf{Correction (post-review, Reviewer AX): the specific ``2 flip points'' claim below is
NOT reproducible and is withdrawn.} My original run found apparent branch disagreements at $N=291,
299$; Reviewer AX's independent reimplementation of the identical continuity-tracking method found
disagreements instead at $N=283,287,289,301,303$ (5 points, not 2), with no disagreement at all at
$N=291$ or $299$. Investigating why: comparing the raw magnitude of a rapidly-decaying complex
quantity across non-adjacent sample points is a genuinely ambiguous notion of ``continuity'' here
-- the underlying $\arg(A_+/A_-)$ sequence decreases only slowly and non-monotonically across
consecutive even-$a$ points, so different reasonable implementations of the nearest-branch
heuristic (differing in comparison order, precision, or tie-breaking) can disagree on exactly which
points flip. \textbf{No specific flip-point claim, or the associated ``real bug in
\texttt{ratio\_amp}'' framing, is retained.} What IS retained and independently confirmed twice
(my own run and Reviewer AX's independent reimplementation, cross-checked against the actual
compiled \texttt{ratio\_amp} binary): the oscillation between the two clusters is genuinely present
under continuity tracking by any of the tested conventions, so it is not an artifact of the naive
per-point principal branch specifically. Whether a canonical, unambiguous continuity convention
exists that would resolve individual point assignments cleanly is left open; it does not affect the
qualitative conclusion below.

\section*{Verdict}
\textbf{Room 7's and Room 9's core findings STAND, un-retracted -- independently confirmed twice
(this note's own computation, and Reviewer AX's separate reimplementation cross-checked against the
compiled binary).} The concern that motivated this check -- that the whole even-$a$ erratic/bimodal
phenomenon might dissolve under a mathematically correct continuous branch choice -- is answered:
\textbf{it does not dissolve.} The phenomenon is real under every continuity convention tested.
The secondary claim of two specific misclassified points is withdrawn (see correction above);
whether \texttt{ratio\_amp} has any real per-point branch bug, and if so which points it affects,
is left genuinely open -- a future fix would need a precisely specified, unambiguous continuity
convention, not the ad hoc nearest-branch heuristic used here.

\paragraph{Weakest points.} Checked only $q=29$, only 12 points, only even-$a$. The two flip
points found here are a small sample; a systematic re-scan of Room 7's full point set with
continuity tracking (not done here, would need porting this logic into \texttt{ratio\_amp} itself
in Rust for the larger $N$ Room 7 actually used) would be the natural next step if this thread is
revisited, but is not needed to answer the yes/no question this note set out to resolve.

\end{document}
